\section{Structure \texttt{button\_config}}

La structure \texttt{button\_config} permet de configurer les boutons GTK4. Elle définit le label, l'icône, le style et le comportement callback du bouton.

\subsection{Définition}

\begin{lstlisting}[language=C]
typedef struct {
    const char *label;
    const char *icon_name;
    const char *theme_variant;
    bool use_underline;
    bool has_frame;
    bool can_shrink;
    int icon_size;
    widget_style_config style;
    void (*on_clicked)(GtkButton *button, gpointer user_data);
    gpointer user_data;
} button_config;

typedef struct {
    const char *label;
    GtkWidget *group_with;
    bool is_active;
    widget_style_config style;
    void (*on_toggled)(GtkCheckButton *button, gpointer user_data);
    gpointer user_data;
} radio_button_config;
\end{lstlisting}

\subsection{Description des champs}

\begin{center}
\begin{tabular}{|l|l|p{5cm}|}
\hline
\textbf{Champ} & \textbf{Type} & \textbf{Description} \\
\hline
label & const char* & Texte affiché sur le bouton \\
\hline
icon\_name & const char* & Icône du thème \\
\hline
theme\_variant & const char* & Variante de thème CSS \\
\hline
use\_underline & bool & Active les mnémoniques \\
\hline
has\_frame & bool & Affiche la bordure \\
\hline
can\_shrink & bool & Autorise la réduction \\
\hline
icon\_size & int & Taille de l'icône \\
\hline
style & widget\_style\_config & Configuration du style \\
\hline
on\_clicked & callback & Fonction appelée au clic \\
\hline
user\_data & gpointer & Données utilisateur \\
\hline
\end{tabular}
\end{center}

\bigskip

Le champ \textbf{label} définit le texte affiché sur le bouton. Peut contenir des mnémoniques ("\_Fichier" pour Alt+F). NULL pour un bouton purement iconique.

Le champ \textbf{icon\_name} spécifie une icône du thème système ("document-open", "edit-copy", etc.). Peut être combiné avec le label.

Le champ \textbf{theme\_variant} applique une variante : "primary" (principal), "success" (positif), "danger" (destructif), "secondary" (secondaire).

Le champ \textbf{use\_underline} active le traitement des mnémoniques clavier.

Le champ \textbf{has\_frame} contrôle l'affichage de la bordure. FALSE crée un bouton plat.

Le champ \textbf{can\_shrink} permet au bouton de réduire sa taille minimale.

Le champ \textbf{icon\_size} définit la taille : GTK\_ICON\_SIZE\_BUTTON (24x24), GTK\_ICON\_SIZE\_SMALL\_TOOLBAR (16x16), etc.

Le champ \textbf{style} est une structure \texttt{widget\_style\_config} pour le style.

Le champ \textbf{on\_clicked} est le callback appelé lors du clic.

Le champ \textbf{user\_data} transmet des données au callback.

\subsection{Boutons radio}

La structure \texttt{radio\_button\_config} permet de créer des boutons radio pour les selections exclusives.

Le champ \textbf{group\_with} lie le bouton à un groupe existant.

Le champ \textbf{is\_active} définit l'état sélectionné initial.

Le champ \textbf{on\_toggled} est appelé lors du changement d'état.

\subsection{Exemple d'utilisation}

\begin{lstlisting}[language=C]
#include "button.h"
#include <stdio.h>

void on_save_clicked(GtkButton *button, gpointer user_data) {
    g_print("Bouton clique!\n");
}

void create_ui_example(GtkWidget *box) {
    button_config save_btn = {
        .label = "_Sauvegarder",
        .theme_variant = "primary",
        .has_frame = true,
        .on_clicked = on_save_clicked,
        .user_data = NULL
    };
    GtkWidget *save_button = create_button(&save_btn);
    gtk_box_append(GTK_BOX(box), save_button);

    radio_button_config opt1 = {
        .label = "Option A",
        .is_active = true
    };
    GtkWidget *radio1 = create_radio_button(&opt1);
    gtk_box_append(GTK_BOX(box), radio1);

    radio_button_config opt2 = {
        .label = "Option B",
        .group_with = radio1
    };
    GtkWidget *radio2 = create_radio_button(&opt2);
    gtk_box_append(GTK_BOX(box), radio2);
}
\end{lstlisting}
